\chapter{Обобщённые координаты микромира}
\label{ch:native-coords}

Геометрия носителя (\cref{ch:carrier}) задаёт узлы~$\Lambda$, шаг~$\hl$ и такт~$\htick$.
Закон~$g$ (\cref{ch:evolution}) действует на регистр ячейки~$\hv$.
Ниже фиксируется \emph{словарь}~$M$: в каких координатах пишется ядро теории
и где заканчивается микромир и начинается мост к~$T$ (CODATA, радианы, $\pi$).
Это не новая решётка и не замена FCC --- только смена записи, как в аналитической механике
при переходе от декартовых координат к обобщённым $(q,\dot q)$.

\begin{remark}[Фаза: $\varphi$, $\omega$, $\Phi$]
\label{rem:phi-omega-Phi}
\textbf{$\varphi_{\mathrm{disc}}\in\mathbb{Z}_{N_{\mathrm{ring}}}$} --- дискретная фаза на кольце (аддитивный индекс $0,\ldots,N_{\mathrm{ring}}-1$).
\textbf{$\omega=e^{i2\pi/N_{\mathrm{ring}}}$} --- генератор вращения (\cref{ch:evolution}); состояние на $S^1$ пишут \emph{мультипликативно} $\omega^{\varphi_{\mathrm{disc}}}$, но координата --- $\varphi_{\mathrm{disc}}$, не $\omega$.
\textbf{$\Phi\in\mathbb{Z}_{N_{\mathrm{ring}}}$} --- фазовый \emph{kick} за такт ($R(\Phi)=\omega^{\Phi}$), не путать с регистром $\varphi_{\mathrm{disc}}$.
В мосте~$T$: $\varphi=\varphi_{\mathrm{disc}}\cdot(2\pi/N_{\mathrm{ring}})$ [рад]; частота $\omega_{\mathrm{SI}}=2\pi/\htick$ --- другое $\omega$.
\end{remark}

\section{Внешние координаты \texorpdfstring{$(dh,dt)$}{(dh,dt)}}
\label{sec:coords-dh-dt}

\begin{definition}[Внешние обобщённые координаты]
\label{def:dh-dt}
На слое~$M$ вводятся счётчики приращений
\[
  dh := \hl := \ell_P,
  \qquad
  dt := \htick.
\]
Позиция узла --- целое число скачков~$dh$ по~$\Lambda$;
время --- счётчик тактов~$dt$.
В натуральных единицах симуляции $dh=dt=1$, $\czero=dh/dt=1$.
\end{definition}

Конфигурация одного узла с внутренним регистром (детали --- \cref{sec:coords-internal}):
\[
  x\in\Lambda,\quad t\in\mathbb{Z},\quad
  Z\in\mathbb{Z}_{N_{\mathrm{ring}}}[i],\quad
  b\in\{0,1\}.
\]

Одна ступень действия~$s_0$ на ячейке за такт порождает механическую лестницу
(\cref{sec:theorem-51}) без трёх независимых подгоняемых параметров:

\begin{table}[ht]
\centering
\caption{Лестница $s_0$ в координатах $(dh,dt)$ и мост к~$T$.}
\label{tab:dh-dt-bridge}
\begin{tabular}{@{}lll@{}}
\toprule
На $M$ & В $(dh,dt)$ & Мост $T$/СИ \\
\midrule
$s_0$ & единица на $(\hv,1\,dt)$ & $\hbar/2$ \\
$p_0$ & $s_0/dh$ & $\hbar/(2\hl)$ \\
$E_0$ & $s_0/dt$ & $\hbar/(2\htick)$ \\
$L_0$ & $s_0$ & $\hbar/2$ \\
$\czero$ & $dh/dt$ & $\hl/\htick$ \\
$F_0$ & $s_0/(dh\cdot dt)$ & $m_{\mathrm{arg}}\,g_M$ \\
$\nu_0$ & $1/dt$ & $1/\htick$ \\
\bottomrule
\end{tabular}
\end{table}

Полный оборот фазового кольца за один такт (\cref{sec:bit-budget}):
\[
  1\,dt \;\Longleftrightarrow\; \text{оборот }\mathbb{Z}_{N_{\mathrm{ring}}},
  \qquad
  N_{\mathrm{ring}}=512.
\]
Один шаг по кольцу --- $\varphi_{\mathrm{disc}}\mapsto\varphi_{\mathrm{disc}}+1\pmod{N_{\mathrm{ring}}}$;
Heisenberg-пол $\Delta\phi_{\mathrm{disc}}=41$ шагов соответствует
$\Delta\phi_{\min}=\tfrac12$ в мосте~$T$.

\begin{remark}[Без $\pi$ в ядре]
На $M$ фазовый остаток пустой ячейки записывается тиками:
\[
  \alpha_*-1=\frac{\Delta\phi_{\mathrm{disc}}}{N_{\mathrm{ring}}}
  =\frac{41}{512};
\]
половина кольца (Паули) --- $N_{\mathrm{ring}}/2=256$ шагов $\varphi_{\mathrm{disc}}$.
Величина $\alpha_*-1\approx 1/(4\pi)$ --- \emph{имя} в мосте~$T$, не определение на~$M$.
\end{remark}

Информационный бюджет ячейки в нативных битах (\cref{sec:bit-budget}):
\[
  B_{\hv}=\log_2 N_{\mathrm{ring}}+\varepsilon_B,
  \qquad
  \varepsilon_B=\{B_{\hv}\}\approx 0{,}065,
  \qquad
  N_{\mathrm{hier}}=\lfloor B_{\hv}\rfloor-1=8.
\]
Пороги CR-энергии (\cref{ch:evolution}) в $(dh,dt)$ при $\czero=1$:
\[
  E_{\mathrm{CR,seed}}=\nu_{\mathrm{CA}}\,\varepsilon_B^2,
  \qquad
  E_{\mathrm{CR,stab}}=\nu_{\mathrm{CA}}\,(1+\varepsilon_B),
\]
где $\nu_{\mathrm{CA}}=\kappa_{\mathrm{link}}\,dh^2/dt$.

\section{Внутренние координаты ячейки \texorpdfstring{$(k_\phi,\varphi_{\mathrm{f}},\mu)$}{(kphi,phif,mu)}}
\label{sec:coords-internal}

Регистр $Z=U+iV\in\mathbb{Z}_{N_{\mathrm{ring}}}[i]$ несёт фазу и амплитуду спинора~$\mathbb{C}^2$.
Число $N_\phi=13=\lceil 4\pi\rceil$ на~$M$ --- целый секторный индекс ($N_{12}+1$),
а не радианная дробь; $13$ не делит $512$, поэтому \emph{мультипликативная}
факторизация $2^9\times 13$ не задаёт естественную решётку.

\begin{definition}[Аддитивные внутренние координаты]
\label{def:internal-coords}
Дискретная фаза $\varphi_{\mathrm{disc}}\in\mathbb{Z}_{N_{\mathrm{ring}}}$ раскладывается по Heisenberg-кванту
$\Delta\phi_{\mathrm{disc}}=41$:
\[
  \varphi_{\mathrm{disc}} \equiv k_\phi\cdot\Delta\phi_{\mathrm{disc}}+\varphi_{\mathrm{f}}
  \pmod{N_{\mathrm{ring}}},
\]
где $k_\phi\in\mathbb{Z}_{N_\phi}$ --- сектор ($0,\ldots,12$),
$\varphi_{\mathrm{f}}\in\mathbb{Z}_{\Delta\phi_{\mathrm{disc}}}$ --- тонкая фаза ($0,\ldots,40$).
Эквивалентно на $S^1$: $\omega^{\varphi_{\mathrm{disc}}}$ при $\omega=e^{i2\pi/N_{\mathrm{ring}}}$ (\cref{rem:phi-omega-Phi}).
Амплитуда на решётке $Q(2^{-B_{\mathrm{amp}}})$, $B_{\mathrm{amp}}=6$:
$\mu\in\mathbb{Z}_{64}$ --- наименьший ненулевой квант $|z|$ на ячейке.
\end{definition}

\begin{proposition}[Каноническая пара $(k_\phi,\varphi_{\mathrm{f}})$]
\label{prop:canonical-phase}
Для каждого $\varphi_{\mathrm{disc}}\in\mathbb{Z}_{512}$ существует единственная каноническая пара
с минимальным $k_\phi$ при $0\le\varphi_{\mathrm{f}}<41$, восстанавливающая $\varphi_{\mathrm{disc}}$.
Наивое отображение $\mathbb{Z}_{13}\times\mathbb{Z}_{41}\to\mathbb{Z}_{512}$
имеет ровно $21$ коллизию; шов (seam) удовлетворяет
\[
  N_{\mathrm{ring}}
  =N_\phi\cdot\Delta\phi_{\mathrm{disc}}-(N_\phi+N_{\mathrm{hier}})
  =13\cdot 41-21.
\]
Пары $(k_\phi=0,\,\varphi_{\mathrm{f}}<21)$ и $(k_\phi=12,\,\varphi_{\mathrm{f}}\ge 21)$
описывают одну и ту же $\varphi_{\mathrm{disc}}$.
\end{proposition}

\begin{proof}
Прямой перебор $512$ значений $\varphi_{\mathrm{disc}}$ и проверка биекции канонической формы
(реализация: \texttt{Internal\_phase\_coords}).
Тождество для шва: $13\cdot 41=533=512+21$, $21=13+8=N_\phi+N_{\mathrm{hier}}$.
\end{proof}

\begin{remark}[Секция вакуума]
Начальное состояние океана (\cref{ch:matter}) сидит на вложении
$\varphi_{\mathrm{f}}=0$, $\varphi_{\mathrm{disc}}=k_\phi\cdot 41$ --- тринадцать ориентаций кирпича,
а не произвольная фаза на кольце.
Эволюция~$g$ заполняет весь $\mathbb{Z}_{512}$.
\end{remark}

Полный спинор: две компоненты $j=1,2$, каждая $(\mu_j,k_{\phi,j},\varphi_{\mathrm{f},j})$
или $(\mu_j,\varphi_{\mathrm{disc},j})$; в вакууме $k_{\phi,1}=k_{\phi,2}$, $\varphi_{\mathrm{f},1}=\varphi_{\mathrm{f},2}=0$.
Относительная фаза SU(2) --- $\varphi_{\mathrm{disc},1}-\varphi_{\mathrm{disc},2}$; динамика на шве из $21$ шага остаётся открытой
(CL-O4, \cref{sec:congruence}).

\section{Мост \texorpdfstring{$M\to T$}{M to T}}
\label{sec:coords-bridge}

\begin{table}[ht]
\centering
\caption{Где живут $\pi$, $\ln 2$ и CODATA.}
\label{tab:m-t-bridge}
\begin{tabular}{@{}p{0.42\linewidth}p{0.48\linewidth}@{}}
\toprule
Ядро $M$ & Только мост $T$ \\
\midrule
$\varphi_{\mathrm{disc}}$, $k_\phi$, $\varphi_{\mathrm{f}}$, $\mu$; $dh$, $dt$; $s_0$ & $\hbar$, $t_P$, $\ell_P$, $c$ \\
$\omega^{\varphi_{\mathrm{disc}}}$, $R(\Phi)=\omega^{\Phi}$ & $e^{i\varphi}$, kick $\Phi$ [рад] \\
$\Delta\phi_{\mathrm{disc}}/N_{\mathrm{ring}}$ & $\alpha_*-1=1/(4\pi)$ \\
$N_{\mathrm{ring}}/2$ шагов $\varphi_{\mathrm{disc}}$ & угол $\pi$ \\
$\log_2 N_{\mathrm{ring}}+\varepsilon_B$ & $B_{\hv}=2\pi/\ln 2$ \\
$N_\phi=13$ & $N_\phi=\lceil 4\pi\rceil$ \\
$k=\Delta x/dh$, $\omega_{\mathrm{SI}}=2\pi/\htick$ & волновой вектор и частота в СИ \\
\bottomrule
\end{tabular}
\end{table}

\begin{remark}
Остаток $\varepsilon_B\approx\{B_{\hv}\}$ --- хвост при записи ёмкости в двоичных координатах
(CR, \cref{ch:evolution}); шов из $21$ шага $\varphi_{\mathrm{disc}}$ --- геометрия фазовой фибрации
(\cref{prop:canonical-phase}). Это разные объекты.
\end{remark}

Далее формулы в главах об эволюции, материи и $\alpha$ можно читать в двух колонках
\cref{tab:dh-dt-bridge,tab:m-t-bridge}: слева --- закон на решётке, справа --- имя наблюдателя.
